\honest{Doublethink (Choosing to be Biased)}{Eliezer Yudkowsky}{2007}

I open with the scene in \textsc{Nineteen Eighty-Four} in which O'Brien burns the photograph and says he does not remember it. Winston fears that ``that lunatic dislocation in the mind could really happen.''

Then I ask: what if self-deception makes us happy? Surely true wisdom would be to choose which biases should govern you. I do not say who holds this view. \nb{There is research behind the question. Taylor and Brown (1988) argued that overly positive views of oneself, of one's control and of the future are normal and ``appear to promote'' mental health, including ``the ability to be happy or contented.'' Their thesis was contested.} I doubt that such a dislocation could really happen.

My first reason is that ``You cannot make yourself believe the sky is green by an act of will.'' \nb{Pascal's advice to the unbeliever was indirect: do what believers did at first, acting ``as if they believed,'' and this ``will naturally make you believe.'' The post does not discuss that route. Two years later, in ``Don't Believe You'll Self-Deceive,'' Yudkowsky wrote of this sentence: ``I wasn't just a dispassionate reporter of the existing facts. I was also trying to instill a self-fulfilling prophecy.''}

My second reason is that a rational choice to be biased needs an accurate model of the consequences first. Then you would have to forget it, and forget the forgetting, as Orwell described. So ``You can't know the consequences of being biased, until you have already debiased yourself,'' and the only alternative is to stay biased ``without any clear idea of the consequences,'' which I call ``willful stupidity.'' \nb{Knowledge of consequences can also come from studies of other people. It is partial, but so is the knowledge behind any decision.}

An optimistic driver skips the seatbelt and is happy for years, until the crash. \nb{Taylor and Brown listed people's belief that they drive better than others among the positive illusions, and suggested that such illusions may help most when a person is threatened or receives bad news. The post's example is a case where the bias leads to a dangerous act.} People leap blindly because dangers they do not know of do not come to mind. ``Been there. Tried that. Got burned.''

A friend told me that the happiness of stupidity is ``really not'' overrated. Maybe stupid people are happier, I grant. ``I suspect'' most of my readers could not reach that happiness, and in the next sentence ``You can never achieve that degree of ignorance.'' I give no evidence for either. The way is closed, and the happiness a rationalist can achieve ``may prove greater, in the end.''
